\documentclass[twocolumn, 12pt]{article}

\usepackage[utf8]{inputenc}
\usepackage[english, spanish]{babel}
\usepackage{fullpage} % changes the margin
\usepackage{graphicx}
\usepackage{amsmath}
\usepackage{enumitem}
\usepackage{chngcntr}
\usepackage{setspace}
\usepackage{url}
\usepackage{csquotes}
\usepackage{float}
\usepackage{verbatim}
\usepackage{tabularx}
\usepackage{amsmath}

\counterwithin{figure}{section}
\renewcommand{\thesection}{\arabic{section}}
\renewcommand{\thesubsection}{\thesection.\arabic{subsection}}
\renewcommand{\baselinestretch}{1.5}

\usepackage[style=apa, maxnames=6, minnames=3, backend=biber]{biblatex}
\DefineBibliographyStrings{english}{%chktex-file 1 chktex-file 6
	andothers = {\em et\addabbrvspace al\adddot}
}
\addbibresource{./Bibliography/bibliography.bib}

\usepackage{array}
\usepackage{enumitem}

\raggedbottom{}

\begin{document}

\begin{titlepage}
	\centering
	\includegraphics[width=0.3\textwidth]{Images/logo_utb.png}\par\vspace{1cm}
	{\scshape\LARGE Universidad Tecnológica de Bolívar \par}
	\vspace{1cm}

	{\scshape\Large FÍSICA ELÉCTRICA \par}
	\vspace{.2cm}

	% chktex-file 8
	{\scshape\Large H1 - C \par}
	\vspace{1cm}
	% chktex-file 8
	\slshape {\Large \bfseries{}Informe de Laboratorio No. V\\}
	\vspace{1cm}

	\slshape {\itshape{} Mauro González, T00067622 \\}
	\slshape {\itshape{} German De Armas Castaño, T00068765 \\}
	\slshape {\itshape{} Angel Vega Rodriguez, T00068186 \\}
	\slshape {\itshape{} Juan Jose Osorio Ariza, T00067316 \\}
	\slshape {\itshape{} Juan Eduardo barón, T00065901 \\}
	\vfill
	Revisado Por \\
	Gabriel Hoyos Gomez Casseres\\
	{\large \today\par}
\end{titlepage}

% ----------------------------------------------------------------------|>
\section{Introducción}

Una bobina es un dispositivo compuesto por un conductor
eléctrico enrollado en forma de espiral o solenoide, que
permite la generación de un campo magnético cuando se le
aplica una corriente eléctrica. El estudio del campo
magnético en una bobina es fundamental en la comprensión de
numerosos fenómenos, tales como la inducción
electromagnética, el funcionamiento de los motores
eléctricos y la generación de energía eléctrica. En este
trabajo, se abordará el análisis del campo magnético en una
bobina desde diferentes perspectivas, incluyendo su origen
y características, la relación entre el campo magnético y
la corriente eléctrica, la influencia de las propiedades
físicas de la bobina en el campo magnético generado y las
aplicaciones prácticas de este fenómeno.

% \newpage

% ----------------------------------------------------------------------|>
\section{Objetivos}

\subsection{Objetivo general}

Calcule el campo magnético producido por el solenoide de CC
debido a la fuerza magnética que actúa sobre la espira
conductora.

\subsection{Objetivo especifico}

\begin{itemize}[label=$\triangleright$]
	\item Comprender la importancia de la estructura del solenoide.
	\item Determinar la importancia de un interruptor para mantener
	      equilibrio en el sistema.
	\item Medir la fuerza magnética ejercida sobre el espiral para
	      mantenerlo en equilibrio.
	\item Comparar la fuerza experimental teórica con la fuerza
	      calculada de forma experimental
\end{itemize}

% ----------------------------------------------------------------------|>
\section{Marco Teórico}

% -------------------------------------------------|>
\subsection*{Campo Magnético}

Es un campo vectorial que describe la interacción magnética
que ocurre entre cargas eléctricas en movimiento. Se
manifiesta como una fuerza que actúa sobre una carga en
movimiento, que se desvía de su trayectoria original en
presencia del campo magnético.

El campo magnético se representa mediante líneas de fuerza
que indican la dirección y la intensidad del campo en cada
punto del espacio. Estas líneas de fuerza forman un patrón
que depende de la distribución de cargas eléctricas en
movimiento y su velocidad. El campo magnético se mide en
unidades de Tesla
\textit{\textbf{(T)}}~(\cite{magnetic-fields-khanacademy}).

% -------------------------------------------------|>
\subsection*{Inducción electromagnética}

La inducción electromagnética se basa en la ley de
Faraday-Neumann-Lenz, que establece que la fuerza
electromotriz (FEM) inducida en un circuito cerrado es
igual al cambio temporal del flujo magnético que atraviesa
el circuito. El flujo magnético es el producto del campo
magnético por la superficie que atraviesa.

% -------------------------------------------------|>
\subsection*{Inducción electromagnética}

La inducción electromagnética se basa en la ley de
Faraday-Neumann-Lenz, que establece que la fuerza
electromotriz (FEM) inducida en un circuito cerrado es
igual al cambio temporal del flujo magnético que atraviesa
el circuito. El flujo magnético es el producto del campo
magnético por la superficie que atraviesa.

% -------------------------------------------------|>
\subsection*{Ley de Biot-Savart}

Esta ley establece que el campo magnético creado por un
conductor eléctrico es directamente proporcional a la
intensidad de la corriente que circula por él y al número
de vueltas de la espira o solenoide. Relaciona los campos
magnéticos con las corrientes que los crean. De una manera
similar a como la ley de Coulomb relaciona los campos
eléctricos con las cargas puntuales que las crean. La
obtención del campo magnético resultante de una
distribución de corrientes implica un producto vectorial, y
cuando la distancia desde la corriente al punto del campo
está variando continuamente, se convierte inherentemente en
un problema de cálculo
diferencial~(\textcite{biot-savart-law}).

% -------------------------------------------------|>
\subsection*{La Regla de la Mano Derecha}

Siempre que circula una corriente eléctrica en un hilo
conductor aparece un campo magnético, cuyas líneas son
circunferencias situadas en planos perpendiculares al
conductor y con sus centros en él. Situamos mentalmente un
sacacorchos de manera que avance en la dirección del polo
positivo al negativo de la pila, los extremos del mango
siguen la dirección de las líneas del campo. A veces la Ley
de Ampère se enuncia empleando el sentido de las agujas del
reloj. Trata de describir el enunciado.

% -------------------------------------------------|>
\subsection*{Formula para calcular el campo magnético dentro de un solenoide}

% chktex-file 24
\begin{equation}
	B = {\frac{\mu_0 N I_b}{L}} \label{eq:FormulaCampoMagneticoSolenoide}
\end{equation}

\nocite{Garcia2016}

Donde: \hfill \break{} \textit{{\large $\mu_o$}}, es la
permeabilidad del espacio libre, \textit{{\large $N$}}, es
la cantidad de vueltas del alambre, \textit{{\large
			$I_b$}}, es la corriente suministrada, \textit{{\large
			$L$}}, la longitud del solenoide.

Dando como resultado (\textit{{\large $B$}}), siendo el
campo magnético del solenoide.

% -------------------------------------------------|>
\subsection*{Formula de fuerza magnética}

\begin{equation}
	Fm = B I_e d \label{eq:FormulaFuerzaMagnetica}
\end{equation}

% ----------------------------------------------------------------------|>
\section{Montaje Experimental}

% ----------------------------------------------------------------------|>
\section{Datos Experimentales}

% --------------------------------------------------|>
% ! Tabla de la toma de datos (1)

% chktex-file 44

\begin{center}
	\begin{tabularx}{0.9\linewidth}{|>{\centering\arraybackslash}X|>{\centering\arraybackslash}X|}
		\hline
		$Fm$ \textit{(mN)} & $I_e$ \textit{(A)} \\ \hline
		0.04               & 0.98               \\ \hline
		0.13               & 1.95               \\ \hline
		0.17               & 3.00               \\ \hline
		0.23               & 3.90               \\ \hline
		0.30               & 5.16               \\ \hline
	\end{tabularx}
\end{center}

\subsection*{Constantes}

\begin{itemize}[label=$\triangleright$]
	\item {\large $N$}: $120$ \textit{(Vueltas)} $\Rightarrow$ {\large $1.20 \times 10^{2}$}
	\item {\large $d$}: $0.04$ \textit{(M)} $\Rightarrow$ {\large $4 \times 10^{-2}$}
	\item {\large $L$}: $0.4$ \textit{(M)} $\Rightarrow$ {\large $4 \times 10^{-1}$}
	\item {\large $\mu_0$}: $4 \pi \times 10^{-7}$ \textit{($T \cdot m/A$)} $\Rightarrow$ {\large $1.26 \times 10^{-6}$}
\end{itemize}

% ----------------------------------------------------------------------|>
\section{Análisis de datos}

% ----------------------------------------------------|>
\subsection*{Calcule el campo magnético en la bobina con la ecuación~(\ref{eq:FormulaCampoMagneticoSolenoide})}

% ! Tabla para el calculo de B Teorico

% chktex-file 44
\begin{tabularx}{0.9\linewidth}{|>{\centering\arraybackslash}X|>{\centering\arraybackslash}X|}
	\hline
	$I_e$ \textit{(A)} & $B$ Teorico \textit{(T)} \textbf{[1]} \\ \hline
	0.98               & $3.69 \times 10^{-4}$                 \\ \hline
	1.95               & $7.35 \times 10^{-4}$                 \\ \hline
	3.00               & $1.13 \times 10^{-3}$                 \\ \hline
	3.90               & $1.47 \times 10^{-3}$                 \\ \hline
	5.16               & $1.95 \times 10^{-3}$                 \\ \hline
\end{tabularx}

\vspace{.5cm}

\begin{tabularx}{0.9\linewidth}{|>{\centering\arraybackslash}X|>{\centering\arraybackslash}X|}
	\hline
	Promedio & $1.13 \times 10^{-3}$ \\ \hline
\end{tabularx}

% ----------------------------------------------------|>
\subsection*{Calcular el campo magnético de la ecuación~(\ref{eq:FormulaFuerzaMagnetica})}

\textit{Despejando {\large $B$} de la ecuación~(\ref{eq:FormulaFuerzaMagnetica})}

\begin{equation}
	B = \frac{Fm}{I_e d} \label{eq:CampoMagneticoDespejado}
\end{equation}

\textit{Tomando la tabla de datos experimentales}

\vspace{.5cm}

% chktex-file 44
\begin{tabularx}{0.9\linewidth}{|>{\centering\arraybackslash}X|}
	\hline
	$B$ Teorico \textit{(T)} \textbf{[2]} \\ \hline
	1.02                                  \\ \hline
	1.67                                  \\ \hline
	1.42                                  \\\hline
	1.47                                  \\ \hline
	1.45                                  \\ \hline
\end{tabularx}

\vspace{.5cm}

\begin{tabularx}{0.9\linewidth}{|>{\centering\arraybackslash}X|>{\centering\arraybackslash}X|}
	\hline
	Promedio & 1.41 \\ \hline
\end{tabularx}

% ----------------------------------------------------|>
\subsection*{Calculando errores}

\begin{tabularx}{0.9\linewidth}{|>{\centering\arraybackslash}X|>{\centering\arraybackslash}X|}
	\hline
	$B$ Teorico \textbf{(1)} & $1.13 \times 10^{-3}$ \\ \hline

	$B$ Teorico \textbf{(2)} & $1.41$                \\ \hline
\end{tabularx}

\vspace{.5cm}

\begin{tabularx}{0.9\linewidth}{|>{\centering\arraybackslash}X|>{\centering\arraybackslash}X|}
	\hline
	Error & 99.91\% \\\hline
\end{tabularx}

Los errores humanos o incluso de los sensores pudieron
haber influido en la toma de datos, resultando en errores
al momento de calcular los campos magnéticos.

% ----------------------------------------------------------------------|>
\section{Conclusiones}

Durante la realización de la experiencia \#7 ``Campo
magnético en una bobina. Fuerza magnética'', se pudo
observar que todo flujo de corriente que pasa a través de
un conductor genera un campo magnético. También se pudo
observar que dos corrientes anti-paralelas, es decir, que
van en direcciones opuestas pero paralelas, experimentan
una fuerza de repulsión magnética entre sí. Por último, se
evidencio que todo campo magnético genera una fuerza con
respecto a un flujo eléctrico o a otra carga. En resumen,
la experiencia permitió comprender la relación entre
corriente y campo magnético, la intensidad eléctrica que
circula por el mismo y la dependencia evidente que existe
entre el campo magnético y la corriente a través de la
bobina.
% \bibliography{./Bibliography/bibliography.bib}
\printbibliography

\end{document}